\documentclass{ocots-td}
\begin{exercise}[label=ex:matrices, points=4, nosolution]
    D'après Sontag, ex. 1.4 : la source est dans le texte d'intro.
    On rappelle les trois axiomes :
    \begin{itemize}
        \item une liste à puces dans l'intro n'est pas une numérotation ;
    \end{itemize}
    \begin{question}
        Montrer que $A$ est inversible. On note
        \[ 1. \quad \text{une formule n'est pas une ligne d'énoncé} \]
        \begin{enumerate}
            \item une liste dans une question énumère des cas ;
            \item elle ne remplace pas une question.
        \end{enumerate}
    \end{question}
    \begin{subquestion}
        En 2024. ce numéro n'est pas en tête de ligne.
    \end{subquestion}
\end{exercise}
\begin{exercise}
    \exercisenotext
    \begin{question}
        Question.
    \end{question}
\end{exercise}
